\documentclass[11pt,a4paper]{article}
\usepackage[margin=27mm]{geometry}
\usepackage{amsmath,amsthm,amssymb,fontspec,unicode-math,hyperref}
\setmainfont{STIXGeneral}
\setmathfont{STIX Two Math}
\setmonofont{Menlo}
\hypersetup{hidelinks}
\usepackage{newunicodechar}
\newfontfamily\numeralfallback{STIX Two Text}
\newunicodechar{∎}{\ensuremath{\blacksquare}}
\newunicodechar{⁰}{{\numeralfallback\char"2070\relax}}
\newunicodechar{₀}{{\numeralfallback\char"2080\relax}}
\newunicodechar{₁}{{\numeralfallback\char"2081\relax}}
\newunicodechar{₂}{{\numeralfallback\char"2082\relax}}
\newunicodechar{₃}{{\numeralfallback\char"2083\relax}}
\newunicodechar{₄}{{\numeralfallback\char"2084\relax}}
\newunicodechar{⁻}{{\numeralfallback\char"207B\relax}}
\newunicodechar{⁸}{{\numeralfallback\char"2078\relax}}
\newunicodechar{ⁱ}{{\numeralfallback\char"2071\relax}}
\newunicodechar{⁴}{{\numeralfallback\char"2074\relax}}
\newunicodechar{⁷}{{\numeralfallback\char"2077\relax}}
\newunicodechar{⁶}{{\numeralfallback\char"2076\relax}}
\newunicodechar{⁵}{{\numeralfallback\char"2075\relax}}
\newunicodechar{₈}{{\numeralfallback\char"2088\relax}}
\newunicodechar{ᵢ}{\textsubscript{i}}
\newunicodechar{ⱼ}{\textsubscript{j}}
\newunicodechar{ₜ}{\textsubscript{t}}
\providecommand{\tightlist}{\setlength{\itemsep}{0pt}\setlength{\parskip}{0pt}}
\newtheorem{theorem}{Theorem}
\newtheorem{proposition}[theorem]{Proposition}
\newtheorem{lemma}[theorem]{Lemma}
\newtheorem{corollary}[theorem]{Corollary}
\theoremstyle{remark}\newtheorem{remark}[theorem]{Remark}
\newcommand{\gen}[1]{\langle #1\rangle}
\newcommand{\normalgen}[1]{\langle\!\langle #1\rangle\!\rangle}
\newcommand{\Z}{\mathbb Z}
\newcommand{\F}{\mathbb F}
\setlength{\emergencystretch}{2em}

\usepackage{longtable,booktabs,array,calc}
\newcounter{none}
\setlength{\emergencystretch}{4em}
\providecommand{\passthrough}[1]{#1}
\title{Central descent and the marked groups of order 1024}
\author{Human Compression Project}
\date{}
\begin{document}
\maketitle
\begin{abstract}
A single finite-kernel parameter space gives a complete nondegenerate
string C-group census at order 1024, with the distinguished marking
retained. We recover all rank, type and marked-facet counts, the 106
finite universal orders and the stated exceptional infinite case.
Central descent gives all three binary lift patterns. Complete raw
presentations identify every named object and expose the literal table
corrections. All 107 source facet-pair rows and all required finite
programs are included; exact algorithms and retained classical
geometric inputs are stated separately.
\end{abstract}
\section{Markings, scope and the main census}
We recover the nondegenerate principal classification in Gomi, Loyola
and De Las Pe\~nas \cite{GLP}. A marking is an ordered generating list
of nonidentity involutions; its parabolics are the subgroups generated
by subsets of that list. It has the intersection property (IP) if the
intersection of any two parabolics is the parabolic on the intersection
of their index sets. A string C-group also has commuting generators at
distance greater than one. Nondegenerate means that each adjacent
product has order greater than two.

The source's equivalence carries the distinguished generator set onto
the other one. For a connected string it therefore identifies an
ordered marking only with itself or its reversal. We retain the marking
through every quotient, presentation and isomorphism test. Abstract
isomorphism of the unmarked groups is not this equivalence.

\begin{theorem}[The principal order-1024 census]
Up to the source's equivalence there are exactly one, 125, 105 and eight
nondegenerate string C-groups of order $1024$ in ranks two, three, four
and five respectively, and none in higher rank. Their rank-three type
counts, marked facet pairs and universal orders are exactly the tables
below. The source's binary-lift and universal-presentation constructions
are recovered with the literal corrections specified in the text.
\end{theorem}

We use the common paper \cite{UA}, specifically UA.Phi, UA.Q,
UA.Forest, UA.Nilpotent, UA.Ladder, UA.Search, UA.Affine and
UA.SquareDecomposition. These are complete proofs, not names for
unproved interfaces. The application supplies the central-descent
argument, exact local presentations, square-center calculations and
the complete source-to-census identification. All four required files
of the shared computational appendix are included without pruning.
The source's infinite universal decomposition and the quoted geometric
and square-classification inputs are stated precisely where used.
\section{Central descent and complete binary lifts}\label{central-descent-and-complete-binary-lifts}

This branch concerns the principal nondegenerate classification in Gomi--Loyola--De Las Peñas \cite{GLP}, with the marking retained. Its equivalence sends the distinguished generator set onto the other set. For a connected string diagram this permits precisely the two orders related by reversal. It does not identify arbitrary markings of an abstractly isomorphic group. The rank-two case is the unique dihedral marking of order 1024. The three higher-rank totals below are 125, 105 and 8; all stated types, facets and universal orders are part of the conclusion.

The specifically retained prior results in this branch are the standard string intersection criteria quoted as \cite{GLP} Lemmas 2.2--2.4, the universal geometric interpretation quoted as its Theorem 4.2, the finite smooth square-quotient classification quoted as its Lemma 5.1, and the exceptional infinite universal decomposition quoted as its Theorem 5.3. Their relevant precise forms are stated below. Exact finite-group and complete coset-enumeration algorithms in GAP, and NQ's universal nilpotent quotient algorithm, are computational inputs. The Atlas supplies definitions of named marked presentations, all included in the defining-data appendix; neither its completeness nor a SmallGroups catalogue is assumed.

\subsection{Preliminary adapters and exact quotient intersections}\label{preliminary-adapters-and-exact-quotient-intersections}

\textbf{Proposition N1 (endpoint cosets).} A finite string C-group of rank d and adjacent orders \(m_{1},\ldots,m_{d-1}\) has order at least \(2\prod m_{i}\). In particular, a nondegenerate string two-group of order 1024 has rank at most five.

\emph{Proof.} Put \(a=r_{0}\), \(b=r_{1}\) and \(H=\langle r_{1},\ldots,r_{d-1}\rangle\). IP gives \(H\cap \langle a,b\rangle=\langle b\rangle\). If \(H(ab)^{i}=H(ab)^{j}\), then \((ab)^{i-j}\) lies in the intersection of the rotation subgroup with \(\langle b\rangle\), which is trivial. The \(m_{1}\) such cosets, with \(0\le i<m_{1}\), are therefore distinct. Iterating gives the bound. Each \(m_{i}\) is at least four in the stated two-group case, so \(2^{2d-1}\)≤1024, forcing \(d\le 5\). ∎

The smooth Coxeter quotient assertion in source Theorem 2.7 is the homomorphism from the string Coxeter presentation sending each distinguished generator to its marked involution; the specified adjacent orders make it smooth. A disconnected string diagram yields a direct product of its component subgroups: they commute and IP makes each meet the subgroup on the other components trivially. This gives the source's preliminary reduction; a separate census of every degenerate product is not added to the principal contract. Source Lemma 2.8 is the coordinate-character consequence of Burnside's basis theorem: an irredundant generating list of a finite two-group maps to a basis of \(G/\Phi(G)\), and every binary assignment extends to a homomorphism to \(C_{2}\). No infinite-group assertion is made here.

For rank three write \(L=\langle a,b\rangle\) and \(R=\langle b,c\rangle\). The retained rank-three criterion says that these two rank-two marked string C-groups, together with \([a,c]=1\) and \(L\cap R=\langle b\rangle\), give full IP. Precisely, for \(d=4\) or 5, the versions quoted in source Lemmas 2.3--2.4 require the prefix \(H=\langle r_{0},\ldots,r_{d-2}\rangle\) to have IP and each equality \(H\cap \langle r_{j},\ldots,r_{d-1}\rangle=\langle r_{j},\ldots,r_{d-2}\rangle\) for \(j=1,\ldots,d-1\), with the final right side trivial. These are respectively all four and all five source hypotheses. All finite certificates below instead test every subset pair directly.

If \(N\lhd G\) and \(N\cap RL=1\), UA.Q says the quotient facets have exactly their old intersection and both embed; their product orders are unchanged. This is source Lemma 3.1. If \(N\le L\) and the displayed images of a,b still give a rank-two marking on \(L/N\), then \(N\cap R\le \langle b\rangle\) and \(b\notin N\), so R embeds. The last clause of UA.Q gives the same facet intersection and the rank-three criterion gives IP. This is source Lemma 3.2, including its symmetric case. An unrelated marking on the abstract quotient \(L/N\) would not suffice.

\subsection{A safe central quotient and the three complete lift patterns}\label{a-safe-central-quotient-and-the-three-complete-lift-patterns}

Assume throughout this subsection that G is a finite marked string two-group \(\langle a,b,c\rangle\), with \(m=|ab|\ge 4\) and \(n=|bc|\ge 4\). Put x=\((ab)^{m/2}\) and y=\((bc)^{n/2}\).

\textbf{Lemma N2 (central elements in an end product).}

\[
Z(G)\cap RL\subseteq\{1,x,y,yx,ybx\}.
\tag{N1}
\]

If \(yx\) is central, x and y are separately central. The element \(ybx\) is never central.

\emph{Proof.} Every element of RL has a unique expression \(h(ab)^{j}\) with \(h\in R\) and \(0\le j<m\). Equality of two expressions would put a rotation in \(L\cap R=\langle b\rangle\). Conjugation by b negates j, so centrality implies \(j=0\) or \(m/2\). The former case gives an element of \(Z(R)=\{1,y\}\). In the latter case write \(z=hx\); since x commutes with b, h lies in \(C_R(b)=\{1,b,y,yb\}\). The choice \(h=b\) fails because the reflection \(bx\) does not commute with a when \(m\ge 4\). This proves (N1).

If \(yx\) commutes with a, then y does, because x already does; commuting with c similarly makes x commute with c.~Thus both half-rotations are central. Finally \(x,y\in \Phi(G)\), since they are even powers of rotations. If \(z=ybx\) were central, then \(a,z,c\) would have the same Frattini basis as \(a,b,c\) and would generate G. A central z is an involution here: \(z\in C_G(x)\) implies y commutes with x, and \(y,b,x\) are commuting involutions. Thus \(a,z,c\) are three commuting involutions, contradicting \(|ab|\ge 4\). ∎

\textbf{Theorem N3 (central descent with IP).} G has a central order-two subgroup N such that the displayed images of \(a,b,c\) give a rank-three string C-group on \(G/N\). Successive such descents, allowing degenerate intermediate markings, stop at \(C_{2}^{3}\).

\emph{Proof.} If x is central, take \(N=\langle x\rangle\). The marked dihedral quotient \(L/N\) has rotation order \(m/2\ge 2\) and R embeds because \(N\cap R=1\). The \(N\le L\) clause of UA.Q proves the quotient intersection. The same reasoning applies to central y. If neither half-rotation is central, choose any central involution z, which exists in every nontrivial finite two-group. Lemma N2 puts z outside RL; now the injective-ends clause of UA.Q proves IP and preserves both end groups. A degenerate rank-three marking is \(C_{2}\) times a dihedral marking. Unless it is \(C_{2}^{3}\), the central rotation involution of its nonabelian dihedral factor can be factored out while preserving rank three. Orders strictly decrease by two, so this process terminates exactly at \(C_{2}^{3}\). ∎

This proves source Corollary 3.3, Lemmas 3.4--3.5 and Theorem 3.6. The bottom boundary is necessary in source Corollary 3.7: a rank-three order-eight group cannot have a rank-three order-four quotient. Its literal unrestricted wording is not asserted.

\textbf{Theorem N4 (three lift patterns with complete presentations).} Let Q be a finite rank-three marked string C-group with a complete presentation in its involutions \(a_{0},b_{0},c_{0}\), including Coxeter relations and additional relators \(R_{1},\ldots,R_{t}\). Let \(G=\langle a,b,c\rangle\) be a central extension with kernel \(\langle z\rangle\) of order two, mapping those involutions respectively to \(a_{0},b_{0},c_{0}\). Suppose the lifts are involutions, \([a,c]=1\) and G has IP. Exactly three order patterns are possible: the left adjacent order doubles, the right adjacent order doubles, or neither doubles. In each case a complete presentation is obtained by adjoining central z with \(z^{2}=1\), each defining Coxeter word of Q set equal to its actual value 1 or z (in particular \((ab)^{p}\)=z when its quotient order p doubles), and equations \(R_{i}=\)\(z^{\epsilon _i}\), \(\epsilon _{i}\in \{0,1\}\). At least one relation expresses z as a word in \(a,b,c\), so z can be eliminated.

\emph{Proof.} Each adjacent order is unchanged or doubled. If both doubled, z would be the halfway rotation in both L and R and hence in their intersection \(\langle b\rangle\), impossible for a nontrivial dihedral rotation. Every defining word of Q lifts to 1 or z. The indicated presentation P maps onto G and has \(P/\langle z\rangle=Q\); centrality and \(z^{2}=1\) imply \(|P|\le 2|Q|=|G|\), proving that map an isomorphism. When an adjacent order doubles, its half-power is z. If neither does and every additional relator lifted to 1, the generating tuple would define a surjection \(Q\to G\), contradicting the orders. ∎

This is source Theorem 3.8 and Table 1. It classifies successful extensions; it does not assert that every binary lift choice has the required order or IP. Those properties and equality of ordered markings must still be checked.

The three successful lift patterns of source Table 1 are therefore:

{\def\LTcaptype{none} % do not increment counter
\begin{longtable}[]{@{}
  >{\raggedright\arraybackslash}p{(\linewidth - 4\tabcolsep) * \real{0.3333}}
  >{\raggedright\arraybackslash}p{(\linewidth - 4\tabcolsep) * \real{0.3333}}
  >{\raggedright\arraybackslash}p{(\linewidth - 4\tabcolsep) * \real{0.3333}}@{}}
\toprule\noalign{}
\begin{minipage}[b]{\linewidth}\raggedright
Pattern
\end{minipage} & \begin{minipage}[b]{\linewidth}\raggedright
Adjacent orders upstairs
\end{minipage} & \begin{minipage}[b]{\linewidth}\raggedright
Defining lifted rotation relations
\end{minipage} \\
\midrule\noalign{}
\endhead
\bottomrule\noalign{}
\endlastfoot
I & \(2p,q\) & \((ab)^p=z,\ (bc)^q=1\) \\
II & \(p,2q\) & \((ab)^p=1,\ (bc)^q=z\) \\
III & \(p,q\) & \((ab)^p=(bc)^q=1\) \\
\end{longtable}
}

Every row additionally retains the lifted involution and string relations, central \(z^2=1\), and all equations \(R_i=z^{\epsilon_i}\). In row III at least one \(\epsilon_i\) is one; this is not sufficient by itself to guarantee that \(z\ne1\) or IP.

\section{Exhaustive bounded classification and named objects}\label{exhaustive-bounded-classification-and-named-objects}

Apply UA.Search (UA.Nilpotent/UA.Ladder) to the path on \(d\) vertices with \(k=\log_2|G|-(2d-1)\). The census program of the shared computational appendix is a single parameterized program for the four cases in the following table. NQ computes the universal nilpotent quotient of class \(k+2\). The original generators are transported through its epimorphism and pc isomorphism; they are never replaced by an arbitrary generating list. Every normal kernel in the invariant ladder is generated, with only equal subgroups deduplicated. The binary canonical subgroup key is used only after all ambient relative pc orders have been checked to be two.

\begin{center}\small
\setlength{\tabcolsep}{3pt}
\begin{tabular}{rrrrrrrrr}
\toprule
$d$ & $k$ & $|U|$ & $|\gamma_3(U)|$ & Kernels & IP & Fixed & Classes & Meets\\
\midrule
3&3&$2^{18}$&$2^{13}$&61&37&3&20&3904\\
3&5&$2^{38}$&$2^{33}$&401&235&15&125&25664\\
4&3&$2^{34}$&$2^{27}$&1483&199&11&105&379648\\
5&1&$2^{16}$&$2^7$&127&15&1&8&130048\\
\bottomrule\end{tabular}\end{center}

Here Kernels counts terminal normal subgroups, IP counts ordered markings satisfying every intersection test, Fixed counts reversal-fixed markings, and Meets counts actual subgroup-intersection checks. The ladder sizes are respectively (7,23,61), (7,23,61,147,401), (31,267,1483) and (127). These are computed outputs, never candidate filters. All 539264 subset-pair intersections are tested, including every candidate rejected by the kernel-profile criterion, and the two IP decisions agree for every candidate. Burnside's orbit formula gives (ordered+fixed)/2 reversal classes. The four ordered type-\(\{4,8\}\) markings in the first case satisfy the four distinct word rows in the companion high-type application, one each; this is the same finite input used in the eight-family completeness proof. The last three cases prove the classification at order 1024, with the rank bound in Proposition N1 proving that no further nondegenerate rank is omitted.

For rank three the twelve type counts, up to reversal, are:

{\def\LTcaptype{none} % do not increment counter
\begin{longtable}[]{@{}lrlr@{}}
\toprule\noalign{}
Type & Count & Type & Count \\
\midrule\noalign{}
\endhead
\bottomrule\noalign{}
\endlastfoot
\(\{4,4\}\) & 1 & \(\{8,8\}\) & 22 \\
\(\{4,8\}\) & 8 & \(\{8,16\}\) & 40 \\
\(\{4,16\}\) & 14 & \(\{8,32\}\) & 8 \\
\(\{4,32\}\) & 4 & \(\{8,64\}\) & 4 \\
\(\{4,64\}\) & 2 & \(\{16,16\}\) & 10 \\
\(\{4,128\}\) & 2 & \(\{16,32\}\) & 10 \\
\end{longtable}
}

These are exactly source Table 5. The complete presentations defining all named facets appear as literal mathematical data in the complete defining-data artifact of the shared computational appendix. All 107 rows of Table 6 are also printed in the table appendix below, using the same bound data. For every one of the 199 rank-four and 15 rank-five ordered markings, the presentation checker of the shared computational appendix reconstructs the quotient from its exact new kernel key and proves unique generator-preserving isomorphisms from both named end groups: 428 exact facet identifications. It computes dual names from reversed generators, rather than assuming that equal-looking types or letter suffixes are self-dual. Under reversal a facet pair \((A,B)\) becomes \((\operatorname{dual}(B),\operatorname{dual}(A))\). Counts are compared after this operation for each of the 99 rank-four and eight rank-five rows individually.

The four printed type-\(\{4,16,4\}\) rows in Table 6 have right-hand type \(\{4,16\}\); their ordered vertex-figure markings necessarily have type \(\{16,4\}\). The data use precisely these duals, preserving every letter suffix and count. This qualification was checked on printed page 16; it is not an inferred change to class counts.

For each of the 106 finite rows, the complete universal presentation is the union of the two complete facet presentations, shifted to overlapping generator lists, together with the distant-end commutation relation. Completed Todd--Coxeter enumeration at the trivial subgroup gives 101 orders of 1024 and five orders of 2048. The finite checker also verifies actual IP and both exact facet embeddings. The remaining row is the explicitly retained infinite universal in the subsection on the exceptional infinite row; it is never sent to an unbounded finite-order test.

\subsection{Every printed rank-three presentation and its corrections}\label{every-printed-rank-three-presentation-and-its-corrections}

The complete finite-presentation protocol in the presentation checker of the shared computational appendix takes the bare displayed relators, performs complete trivial-subgroup Todd--Coxeter enumeration, and checks that the resulting permutation action has the table's full order. It evaluates every relator, proves generation by the specified marked involutions, and tests all \(4^{d}\) marked-subgroup intersections. Thus its order bounds concern the raw presented groups, not just finite images or nilpotent truncations.

Write \(C(p,q)\) for the rank-three Coxeter relations of type \(\{p,q\}\). In source Table 2, adjoin to \(C(4,8)\) the following pairs, set equal to 1:

{\def\LTcaptype{none} % do not increment counter
\begin{longtable}[]{@{}lll@{}}
\toprule\noalign{}
Name & \(R_{1}\) & \(R_{2}\) \\
\midrule\noalign{}
\endhead
\bottomrule\noalign{}
\endlastfoot
512a & \((abcb)^{4}\) & \((ca(bc)^{3}ab)^{2}\) \\
512b & \((cab)^{8}\) & \((a(bc)^{3}b)^{2}\) \\
512c & \((c(abcb)^{2})^{2}\) & \((ca(bc)^{3}ab)^{2}\) \\
512d & \((a(bc)^{3}b)^{2}\) & \((cab)^{5}c(ba)^{2}(cb)^{2}ab\) \\
\end{longtable}
}

All four raw groups have order 512 and are pairwise inequivalent as ordered markings; their complete primary Atlas presentations are identified exactly. Class I of source Example 3.9 would require an order-512 quotient of \(C(2,8)\), whose order is 32, so contributes none. Class II uses \(C(4,4)\) with \((abcb)^{8}=1\), of order 512; its two lifts have central \(\tau =(bc)^{4}\) and \((abcb)^{8}=1\) or \(\tau\). They define the two distinct targets called 1024a and 1024b. For each Table 2 row, the three Class III choices \((R_{1},R_{2})=(\tau ,1),(1,\tau ),(\tau ,\tau )\), with central \(\tau ^{2}=1\), yield exactly:

{\def\LTcaptype{none} % do not increment counter
\begin{longtable}[]{@{}llll@{}}
\toprule\noalign{}
Base & \((\tau ,1)\) & \((1,\tau )\) & \((\tau ,\tau )\) \\
\midrule\noalign{}
\endhead
\bottomrule\noalign{}
\endlastfoot
512a & 1024c & 1024d & 1024e \\
512b & 1024a & 1024c & 1024f \\
512c & 1024c & 1024g & 1024h \\
512d & 512d & 1024a & 512d \\
\end{longtable}
}

Quotienting by the actual \(\tau\) subgroup recovers the named base marking in every case. In the two order-512 outcomes \(\tau =1\). All eight order-1024 targets are pairwise inequivalent as ordered markings and match the eight type-\(\{4,8\}\) quotients in the unfiltered census bijectively, proving completeness independent of any source table assertion.

Separately evaluating every literal Table 4 row gives the ordered-label map

\[
(a,b,c,d,e,f,g,h)_{\mathrm{printed}}
\longmapsto(a,b,c,d,c,a,g,c)_{\mathrm{Table\,3}}.
\tag{N2}
\]

To obtain its intended eight targets, the second displayed additional relation in each of printed rows e,f,h must have right-hand side \(\tau\) instead of 1. These are exactly the already checked \((\tau ,\tau )\) lifts of Table 2 rows \(a,b,c\). No other row or count changes. Printed page 10 was inspected visually; the diagnosis is about the literal source, without claiming priority or unique editorial intent for the repair.

The literal Table 4 presentations may be read without guessing from the names. In each row use the type-\(\{4,8\}\) Coxeter relations and a central generator \(\tau\) with \(\tau^2=1\). Rows a,b impose \(\tau=(bc)^4\); the other rows impose the indicated pair of words from the stated Table 2 base:

{\def\LTcaptype{none} % do not increment counter
\begin{longtable}[]{@{}lll@{}}
\toprule\noalign{}
Printed row & Base or extra relation & Right-hand sides \\
\midrule\noalign{}
\endhead
\bottomrule\noalign{}
\endlastfoot
a & \((abcb)^8\), with \(\tau=(bc)^4\) & \(1\) \\
b & \((abcb)^8\), with \(\tau=(bc)^4\) & \(\tau\) \\
c & Table 2 base a: \((R_1,R_2)\) & \((\tau,1)\) \\
d & Table 2 base a: \((R_1,R_2)\) & \((1,\tau)\) \\
e & Table 2 base a: \((R_1,R_2)\) & \((\tau,1)\) \\
f & Table 2 base b: \((R_1,R_2)\) & \((\tau,1)\) \\
g & Table 2 base c: \((R_1,R_2)\) & \((1,\tau)\) \\
h & Table 2 base c: \((R_1,R_2)\) & \((\tau,1)\) \\
\end{longtable}
}

Thus the literal rows and their corrected intended targets remain distinct records. The correction does not silently redefine the raw presentations which the checker evaluates.

\subsection{Universal presentations and all finite examples}\label{universal-presentations-and-all-finite-examples}

For source Definition 4.1 use the combined facet presentation described above. The retained geometric universal principle, source Theorem 4.2, states that a nonempty class of regular polytopes with those specified facet and vertex-figure groups has that marked universal string C-group and every member is its smooth quotient. At the group-presentation level the epimorphism is immediate from the defining words; the abstract-polytope interpretation remains an explicit background input.

The bare presentations in Examples 4.3--4.5 give:

{\def\LTcaptype{none} % do not increment counter
\begin{longtable}[]{@{}
  >{\raggedright\arraybackslash}p{(\linewidth - 6\tabcolsep) * \real{0.2308}}
  >{\raggedleft\arraybackslash}p{(\linewidth - 6\tabcolsep) * \real{0.3077}}
  >{\raggedright\arraybackslash}p{(\linewidth - 6\tabcolsep) * \real{0.2308}}
  >{\raggedright\arraybackslash}p{(\linewidth - 6\tabcolsep) * \real{0.2308}}@{}}
\toprule\noalign{}
\begin{minipage}[b]{\linewidth}\raggedright
Object
\end{minipage} & \begin{minipage}[b]{\linewidth}\raggedleft
Raw order
\end{minipage} & \begin{minipage}[b]{\linewidth}\raggedright
Type
\end{minipage} & \begin{minipage}[b]{\linewidth}\raggedright
Ordered end labels
\end{minipage} \\
\midrule\noalign{}
\endhead
\bottomrule\noalign{}
\endlastfoot
Example 4.3 universal & 512 & \(\{4,4,4\}\) & \(\{4,4\}\)*32, \(\{4,4\}\)*128 \\
Example 4.4 universal & 1024 & \(\{4,4,4\}\) & \(\{4,4\}\)*32, \(\{4,4\}\)*256 \\
Example 4.5 universal & 2048 & \(\{4,4,8\}\) & \(\{4,4\}\)*64, \(\{4,8\}\)*128a \\
Example 4.5 first quotient & 1024 & \(\{4,4,8\}\) & same two labels \\
Example 4.5 second quotient & 1024 & \(\{4,4,8\}\) & same two labels \\
Example 4.5 third quotient & 1024 & \(\{4,4,4\}\) & not in the prescribed smooth class \\
\end{longtable}
}

The first two universal presentations add \((abcb)^{2}=(bcdc)^{4}=1\) and \((abcb)^{2}=(dbc)^{8}=1\) respectively to the type-\(\{4,4,4\}\) Coxeter relations. The third adds

\[
(cab)^4=1,\qquad dbcd(cb)^2(dc)^2bc=1
\tag{N3}
\]

to the type-\(\{4,4,8\}\) relations. The complete named facet words in the complete defining-data artifact of the shared computational appendix differ, where necessary, only by cyclic conjugation from these displayed source words. Their exact generator-preserving isomorphisms and the full raw order/IP tests verify each row above.

For the order-2048 group the three printed central-involution words are

\[
bcdcbab(dc)^2dba,\qquad(bcdcba)^2,\qquad(cd)^4.
\tag{N4}
\]

The checker independently enumerates all central involutions and obtains exactly these three distinct elements. Adding each as a raw relator gives the three quotient rows above. The first two are inequivalent marked groups, with unchanged type and end markings. Every normal order-two subgroup is central, so these three exhaust all order-1024 quotients of this universal, and exactly two belong to the prescribed smooth class. Example 4.3 has no quotient of order 1024; Example 4.4 has only itself. These statements concern these exact orders and do not extrapolate centrality to larger kernels.

\subsection{Square-lattice centers and the exceptional infinite row}\label{square-lattice-centers-and-the-exceptional-infinite-row}

Retain the finite smooth square-quotient classification quoted in source Lemma 5.1: a finite smooth quotient of \([4,4]\) lies in one of the two families obtained by adding \((abcb)^{s}\)=1, of order \(8s^{2}\), or \((abc)^{2s}\)=1, of order \(16s^{2}\), for a positive integer s. The former has \(|abcb|=s\) and \(|abc|=2s\); the latter has \(|abcb|=2s\) and \(|abc|=2s\). The square and diagonal decompositions are now supplied by UA.SquareDecomposition, including \(s=1\). In the marking used here they are \[
 (\langle a,c^b\rangle\times\langle c,a^b\rangle)\rtimes\langle b\rangle,
 \qquad
 (\langle b,b^{ac}\rangle\times\langle b^a,b^c\rangle)\rtimes\langle a,c\rangle,
\] respectively. Each factor has order \(2s\) and the complements have orders two and four. To check the marking in the second expression, its two pairs of reflection axes are \(x-y=0,1\) and \(x+y=1,0\); their products translate by \((-1,1)\) and \((1,1)\). The perpendicular factors, independent translations and complement parity check are exactly those in UA.SquareDecomposition. The source instead prints the pairs \(\langle b,b^a\rangle\), \(\langle b^c,b^{ac}\rangle\). That pairing is erroneous for general s: \(bb^a=(ba)^2\) has order two, so its first factor always has order four. The classification is retained; this structural correction and its proof are not assumed from it. Only the finite scope is used; the infinite Coxeter group itself is not asserted to admit such an s.

Use the common affine model of UA.Affine after reversing its endpoints and conjugating all three generators by b. Since the lattices are invariant under the coordinate swap, the same square and diagonal quotients have coordinates

\[
a(x,y)=(1-x,y),\quad b(x,y)=(y,x),\quad c(x,y)=(x,-y).
\tag{N5}
\]

Its group is \(W=\mathbb Z^{2}\rtimes D_{8}\); \(t=abcb\) translates by \((1,0)\), \((abc)^{2}\) translates by \((1,1)\), \((bc)^{2}\) is \((0,-I)\), and \((ab)^{2}\) is \(((1,1),-I)\). The two finite families are \(W/s\mathbb Z^{2}\) and \(W/L_s\), where \(L_s=\langle (s,s),(s,-s)\rangle\). The UA.Affine normal forms establish these identifications; equivalently the displayed source relations give surjections onto these affine images of the stated orders, hence isomorphisms.

\textbf{Proposition N5 (all source center formulas).} For s=\(2^{m}\), in \(W/s\mathbb Z^{2}\) the center is \(\langle (ab)^{2},(bc)^{2}\rangle\) when \(m=1\) and ⟨\((abc)^{s}\)⟩ when \(m\ge 2\). In \(W/L_s\) it is \(\langle ac,(ab)^{2},(bc)^{2}\rangle\) when \(m=0\) and ⟨\((abcb)^{s}\)⟩ when \(m\ge 1\).

\emph{Proof.} In a central affine pair \((v,M)\), commutation with all translations forces M to act trivially on \(T=\mathbb Z^{2}/L\). If the signed-permutation action on T is faithful, \(M=I\) and the center is \(T^{D_8}\).

For \(L=s\mathbb Z^{2}\), \(s\ge 4\), faithfulness is immediate on the two coordinate vectors. Invariance gives \(2x=2y=0\) and \(x=y\) modulo s, leaving precisely 0 and \((s/2,s/2)\). For \(s=2\), the action kernel is the diagonal sign subgroup; to commute also with the linear \(D_{8}\) subgroup, M must lie in \(\{I,-I\}\) and v in \({0,(1,1)}\). This gives the first exceptional center. No first-family formula at \(m=0\) is claimed in the source or here.

For \(L=L_s\) with \(s\ge 2\), a nontrivial diagonal sign change cannot act trivially, since \((2,0)\) and \((0,2)\) do not lie in \(L_s\). A signed swap cannot act trivially because its difference on a coordinate vector has a coordinate \(\pm 1\), not divisible by s. A fixed translation satisfies \((2x,0),(0,2y)\in L_s\), hence s divides both \(x,y\); modulo \(L_s\) its class is 0 or \((s,0)\), both swap-invariant. For \(s=1\) the lattice consists of vectors with even coordinate sum, so \(T=C_{2}\) with trivial \(D_{8}\) action and the group is \(C_{2}\times D_{8}\). Here \(ac=(t,-I)\), \((bc)^{2}=(0,-I)\), and \((ab)^{2}=(bc)^{2}\) in the quotient, giving exactly the claimed generating set. Its three displayed generators are not independent. ∎

For the sole infinite row of Table 6 retain the exact source Theorem 5.3, attributed there to McMullen--Schulte: the universal type-\(\{4,4,4\}\) group with extra relations \((abcb)^{4}=(bcd)^{4}=1\) has

\[
U=(W_1\times W_2)\rtimes\langle b,d\rangle,
\quad \langle b,d\rangle\cong C_2^2,\quad
W_1=\langle c^{bd},a,c^b\rangle\cong[4,4],\quad
W_2=\langle c,a^b,c^d\rangle=W_1^b.
\tag{N6}
\]

In particular it is infinite. This theorem, including its marked decomposition, is retained and is not a result of finite coset enumeration.

\textbf{Proposition N6 (the exceptional kernel location).} If \(\pi :U\to G\) is a marked epimorphism to a finite string two-group with IP, then \(\ker\pi\le W_{1}\times W_{2}\) and

\[
G\cong ((W_1\times W_2)/\ker\pi)\rtimes\langle\pi(b),\pi(d)\rangle.
\tag{N7}
\]

\emph{Proof.} IP makes the four marked generators irredundant. Their images form a basis of \(G/\Phi(G)\). The b- and d-coordinate characters, composed with π, agree on every generator with the projection \(U\to C_{2}^{2}\) in (N6): the displayed \(W_{i}\) generators have even b- and d-parities. Thus every element of \(\ker\pi\) lies in \(W_{1}\times W_{2}\). The image of \(\langle b,d\rangle\) still has order four by these same coordinate characters and meets the image of \(W_{1}\times W_{2}\) trivially, proving the split quotient formula. ∎

This proves source Lemma 5.4 without assuming quotients preserve IP. At order 1024 the first factor in (N7) has order 256. Finally the two source Theorem 5.5 presentations share the type-\(\{4,4,4\}\) Coxeter relators and \((abcb)^{4}=(bcd)^{4}=1\). The first adds \((dcbabdc)^{2}=1\). The second adds

\[
(bdcdba)^2(cbab)^2=1,\qquad(abcb)^2(babdcd)^2=1.
\tag{N8}
\]

The presentation checker of the shared computational appendix proves directly that these two raw groups have order 1024 and full IP, and identifies them with the two and only two census markings whose ordered ends are \(\{4,4\}\)*128 and \(\{4,4\}\)*64. Unequal end orders rule out their becoming equal under reversal. This proves the exact exceptional classification without repeating the source's separate case analysis or presuming its presentations finite.

The finite execution in the presentation checker of the shared computational appendix totals 214 complete presentations and 43840 actual IP meets: 72 defining Atlas models, 34 source example/literal-table presentations, the two exceptional presentations, and 106 finite universal presentations. Its distinct logical role from the census program of the shared computational appendix is preserved: the universal-kernel proof gives exhaustion, while complete raw coset tables and exact marked isomorphisms give presentation and name recognition. All code and mathematical data are included as complete external appendices with the bindings in the shared computational appendix. These are exact finite computations conditional on the stated algorithms, not Lean proofs or experimental extrapolations. The concluding source suggestion about geometric realizations, a separate classification of all degenerate products, and historical novelty of the corrections are outside the contract.

\appendix
\section{All marked facet pairs and universal orders}
The following 107 rows reproduce the complete bound mathematical data,
with the four right-hand $\{4,16\}$ labels written as their necessary
ordered $\{16,4\}$ duals. A star in a label separates its type from its
order and suffix; it does not assert any additional duality. The last
eight rows have rank five. $u$ is the universal group order and $c$ is
the number of order-1024 classes. The single $\infty$ entry invokes the
retained marked decomposition; every finite entry is checked by
completed trivial-subgroup coset enumeration. The independent kernel
census identifies the named ends and obtains the class counts; these
printed numbers are comparison outputs, not search filters.

\small
\begin{longtable}{rllllr}
\toprule
Row & Type & First end & Second end & $u$ & $c$\\
\midrule\endfirsthead
\toprule Row & Type & First end & Second end & $u$ & $c$\\\midrule\endhead
\bottomrule\endfoot
1 & \(\{4,4,4\}\) & \(\{4,4\}\!\ast 32\) & \(\{4,4\}\!\ast 256\) & 1024 & 1 \\
2 & \(\{4,4,4\}\) & \(\{4,4\}\!\ast 64\) & \(\{4,4\}\!\ast 64\) & 1024 & 1 \\
3 & \(\{4,4,4\}\) & \(\{4,4\}\!\ast 128\) & \(\{4,4\}\!\ast 64\) & \(\infty\) & 2 \\
4 & \(\{4,4,8\}\) & \(\{4,4\}\!\ast 32\) & \(\{4,8\}\!\ast 256a\) & 1024 & 1 \\
5 & \(\{4,4,8\}\) & \(\{4,4\}\!\ast 32\) & \(\{4,8\}\!\ast 256b\) & 1024 & 1 \\
6 & \(\{4,4,8\}\) & \(\{4,4\}\!\ast 32\) & \(\{4,8\}\!\ast 256c\) & 1024 & 1 \\
7 & \(\{4,4,8\}\) & \(\{4,4\}\!\ast 32\) & \(\{4,8\}\!\ast 256d\) & 1024 & 1 \\
8 & \(\{4,4,8\}\) & \(\{4,4\}\!\ast 64\) & \(\{4,8\}\!\ast 128a\) & 2048 & 2 \\
9 & \(\{4,4,8\}\) & \(\{4,4\}\!\ast 64\) & \(\{4,8\}\!\ast 128b\) & 2048 & 2 \\
10 & \(\{4,4,8\}\) & \(\{4,4\}\!\ast 128\) & \(\{4,8\}\!\ast 64a\) & 1024 & 1 \\
11 & \(\{4,4,8\}\) & \(\{4,4\}\!\ast 128\) & \(\{4,8\}\!\ast 64b\) & 1024 & 1 \\
12 & \(\{4,4,16\}\) & \(\{4,4\}\!\ast 32\) & \(\{4,16\}\!\ast 256a\) & 1024 & 1 \\
13 & \(\{4,4,16\}\) & \(\{4,4\}\!\ast 32\) & \(\{4,16\}\!\ast 256b\) & 1024 & 1 \\
14 & \(\{4,4,16\}\) & \(\{4,4\}\!\ast 64\) & \(\{4,16\}\!\ast 128a\) & 1024 & 1 \\
15 & \(\{4,4,16\}\) & \(\{4,4\}\!\ast 64\) & \(\{4,16\}\!\ast 128b\) & 1024 & 1 \\
16 & \(\{4,4,32\}\) & \(\{4,4\}\!\ast 32\) & \(\{4,32\}\!\ast 256a\) & 1024 & 1 \\
17 & \(\{4,4,32\}\) & \(\{4,4\}\!\ast 32\) & \(\{4,32\}\!\ast 256b\) & 1024 & 1 \\
18 & \(\{4,8,4\}\) & \(\{4,8\}\!\ast 64a\) & \(\{8,4\}\!\ast 256a\) & 1024 & 1 \\
19 & \(\{4,8,4\}\) & \(\{4,8\}\!\ast 64a\) & \(\{8,4\}\!\ast 256b\) & 1024 & 1 \\
20 & \(\{4,8,4\}\) & \(\{4,8\}\!\ast 64a\) & \(\{8,4\}\!\ast 256c\) & 1024 & 1 \\
21 & \(\{4,8,4\}\) & \(\{4,8\}\!\ast 64a\) & \(\{8,4\}\!\ast 256d\) & 1024 & 1 \\
22 & \(\{4,8,4\}\) & \(\{4,8\}\!\ast 64b\) & \(\{8,4\}\!\ast 256a\) & 1024 & 1 \\
23 & \(\{4,8,4\}\) & \(\{4,8\}\!\ast 64b\) & \(\{8,4\}\!\ast 256b\) & 1024 & 1 \\
24 & \(\{4,8,4\}\) & \(\{4,8\}\!\ast 64b\) & \(\{8,4\}\!\ast 256c\) & 1024 & 1 \\
25 & \(\{4,8,4\}\) & \(\{4,8\}\!\ast 64b\) & \(\{8,4\}\!\ast 256d\) & 1024 & 1 \\
26 & \(\{4,8,4\}\) & \(\{4,8\}\!\ast 128a\) & \(\{8,4\}\!\ast 128a\) & 2048 & 2 \\
27 & \(\{4,8,4\}\) & \(\{4,8\}\!\ast 128a\) & \(\{8,4\}\!\ast 128b\) & 2048 & 2 \\
28 & \(\{4,8,4\}\) & \(\{4,8\}\!\ast 128b\) & \(\{8,4\}\!\ast 128b\) & 2048 & 2 \\
29 & \(\{4,8,8\}\) & \(\{4,8\}\!\ast 64a\) & \(\{8,8\}\!\ast 256a\) & 1024 & 1 \\
30 & \(\{4,8,8\}\) & \(\{4,8\}\!\ast 64a\) & \(\{8,8\}\!\ast 256b\) & 1024 & 1 \\
31 & \(\{4,8,8\}\) & \(\{4,8\}\!\ast 64a\) & \(\{8,8\}\!\ast 256c\) & 1024 & 1 \\
32 & \(\{4,8,8\}\) & \(\{4,8\}\!\ast 64a\) & \(\{8,8\}\!\ast 256d\) & 1024 & 1 \\
33 & \(\{4,8,8\}\) & \(\{4,8\}\!\ast 64a\) & \(\{8,8\}\!\ast 256e\) & 1024 & 1 \\
34 & \(\{4,8,8\}\) & \(\{4,8\}\!\ast 64a\) & \(\{8,8\}\!\ast 256f\) & 1024 & 1 \\
35 & \(\{4,8,8\}\) & \(\{4,8\}\!\ast 64a\) & \(\{8,8\}\!\ast 256g\) & 1024 & 1 \\
36 & \(\{4,8,8\}\) & \(\{4,8\}\!\ast 64a\) & \(\{8,8\}\!\ast 256h\) & 1024 & 1 \\
37 & \(\{4,8,8\}\) & \(\{4,8\}\!\ast 64b\) & \(\{8,8\}\!\ast 256a\) & 1024 & 1 \\
38 & \(\{4,8,8\}\) & \(\{4,8\}\!\ast 64b\) & \(\{8,8\}\!\ast 256b\) & 1024 & 1 \\
39 & \(\{4,8,8\}\) & \(\{4,8\}\!\ast 64b\) & \(\{8,8\}\!\ast 256c\) & 1024 & 1 \\
40 & \(\{4,8,8\}\) & \(\{4,8\}\!\ast 64b\) & \(\{8,8\}\!\ast 256d\) & 1024 & 1 \\
41 & \(\{4,8,8\}\) & \(\{4,8\}\!\ast 64b\) & \(\{8,8\}\!\ast 256e\) & 1024 & 1 \\
42 & \(\{4,8,8\}\) & \(\{4,8\}\!\ast 64b\) & \(\{8,8\}\!\ast 256f\) & 1024 & 1 \\
43 & \(\{4,8,8\}\) & \(\{4,8\}\!\ast 64b\) & \(\{8,8\}\!\ast 256g\) & 1024 & 1 \\
44 & \(\{4,8,8\}\) & \(\{4,8\}\!\ast 64b\) & \(\{8,8\}\!\ast 256h\) & 1024 & 1 \\
45 & \(\{4,8,8\}\) & \(\{4,8\}\!\ast 128a\) & \(\{8,8\}\!\ast 128a\) & 1024 & 1 \\
46 & \(\{4,8,8\}\) & \(\{4,8\}\!\ast 128a\) & \(\{8,8\}\!\ast 128b\) & 1024 & 1 \\
47 & \(\{4,8,8\}\) & \(\{4,8\}\!\ast 128a\) & \(\{8,8\}\!\ast 128c\) & 1024 & 1 \\
48 & \(\{4,8,8\}\) & \(\{4,8\}\!\ast 128a\) & \(\{8,8\}\!\ast 128d\) & 1024 & 1 \\
49 & \(\{4,8,8\}\) & \(\{4,8\}\!\ast 128b\) & \(\{8,8\}\!\ast 128a\) & 1024 & 1 \\
50 & \(\{4,8,8\}\) & \(\{4,8\}\!\ast 128b\) & \(\{8,8\}\!\ast 128b\) & 1024 & 1 \\
51 & \(\{4,8,8\}\) & \(\{4,8\}\!\ast 128b\) & \(\{8,8\}\!\ast 128c\) & 1024 & 1 \\
52 & \(\{4,8,8\}\) & \(\{4,8\}\!\ast 128b\) & \(\{8,8\}\!\ast 128d\) & 1024 & 1 \\
53 & \(\{4,8,16\}\) & \(\{4,8\}\!\ast 64a\) & \(\{8,16\}\!\ast 256a\) & 1024 & 1 \\
54 & \(\{4,8,16\}\) & \(\{4,8\}\!\ast 64a\) & \(\{8,16\}\!\ast 256b\) & 1024 & 1 \\
55 & \(\{4,8,16\}\) & \(\{4,8\}\!\ast 64a\) & \(\{8,16\}\!\ast 256c\) & 1024 & 1 \\
56 & \(\{4,8,16\}\) & \(\{4,8\}\!\ast 64a\) & \(\{8,16\}\!\ast 256d\) & 1024 & 1 \\
57 & \(\{4,8,16\}\) & \(\{4,8\}\!\ast 64a\) & \(\{8,16\}\!\ast 256e\) & 1024 & 1 \\
58 & \(\{4,8,16\}\) & \(\{4,8\}\!\ast 64a\) & \(\{8,16\}\!\ast 256f\) & 1024 & 1 \\
59 & \(\{4,8,16\}\) & \(\{4,8\}\!\ast 64b\) & \(\{8,16\}\!\ast 256c\) & 1024 & 1 \\
60 & \(\{4,8,16\}\) & \(\{4,8\}\!\ast 64b\) & \(\{8,16\}\!\ast 256d\) & 1024 & 1 \\
61 & \(\{4,8,16\}\) & \(\{4,8\}\!\ast 64b\) & \(\{8,16\}\!\ast 256e\) & 1024 & 1 \\
62 & \(\{4,8,16\}\) & \(\{4,8\}\!\ast 64b\) & \(\{8,16\}\!\ast 256f\) & 1024 & 1 \\
63 & \(\{4,16,4\}\) & \(\{4,16\}\!\ast 128a\) & \(\{16,4\}\!\ast 256a\) & 1024 & 1 \\
64 & \(\{4,16,4\}\) & \(\{4,16\}\!\ast 128a\) & \(\{16,4\}\!\ast 256b\) & 1024 & 1 \\
65 & \(\{4,16,4\}\) & \(\{4,16\}\!\ast 128b\) & \(\{16,4\}\!\ast 256a\) & 1024 & 1 \\
66 & \(\{4,16,4\}\) & \(\{4,16\}\!\ast 128b\) & \(\{16,4\}\!\ast 256b\) & 1024 & 1 \\
67 & \(\{4,16,8\}\) & \(\{4,16\}\!\ast 128a\) & \(\{16,8\}\!\ast 256a\) & 1024 & 1 \\
68 & \(\{4,16,8\}\) & \(\{4,16\}\!\ast 128a\) & \(\{16,8\}\!\ast 256b\) & 1024 & 1 \\
69 & \(\{4,16,8\}\) & \(\{4,16\}\!\ast 128a\) & \(\{16,8\}\!\ast 256c\) & 1024 & 1 \\
70 & \(\{4,16,8\}\) & \(\{4,16\}\!\ast 128a\) & \(\{16,8\}\!\ast 256d\) & 1024 & 1 \\
71 & \(\{4,16,8\}\) & \(\{4,16\}\!\ast 128a\) & \(\{16,8\}\!\ast 256e\) & 1024 & 1 \\
72 & \(\{4,16,8\}\) & \(\{4,16\}\!\ast 128a\) & \(\{16,8\}\!\ast 256f\) & 1024 & 1 \\
73 & \(\{4,16,8\}\) & \(\{4,16\}\!\ast 128b\) & \(\{16,8\}\!\ast 256a\) & 1024 & 1 \\
74 & \(\{4,16,8\}\) & \(\{4,16\}\!\ast 128b\) & \(\{16,8\}\!\ast 256b\) & 1024 & 1 \\
75 & \(\{4,16,8\}\) & \(\{4,16\}\!\ast 128b\) & \(\{16,8\}\!\ast 256c\) & 1024 & 1 \\
76 & \(\{4,16,8\}\) & \(\{4,16\}\!\ast 128b\) & \(\{16,8\}\!\ast 256d\) & 1024 & 1 \\
77 & \(\{4,16,8\}\) & \(\{4,16\}\!\ast 128b\) & \(\{16,8\}\!\ast 256e\) & 1024 & 1 \\
78 & \(\{4,16,8\}\) & \(\{4,16\}\!\ast 128b\) & \(\{16,8\}\!\ast 256f\) & 1024 & 1 \\
79 & \(\{4,32,4\}\) & \(\{4,32\}\!\ast 256a\) & \(\{32,4\}\!\ast 256a\) & 1024 & 1 \\
80 & \(\{4,32,4\}\) & \(\{4,32\}\!\ast 256a\) & \(\{32,4\}\!\ast 256b\) & 1024 & 1 \\
81 & \(\{4,32,4\}\) & \(\{4,32\}\!\ast 256b\) & \(\{32,4\}\!\ast 256b\) & 1024 & 1 \\
82 & \(\{8,4,8\}\) & \(\{8,4\}\!\ast 64a\) & \(\{4,8\}\!\ast 128a\) & 1024 & 1 \\
83 & \(\{8,4,8\}\) & \(\{8,4\}\!\ast 64a\) & \(\{4,8\}\!\ast 128b\) & 1024 & 1 \\
84 & \(\{8,4,8\}\) & \(\{8,4\}\!\ast 64b\) & \(\{4,8\}\!\ast 128a\) & 1024 & 1 \\
85 & \(\{8,4,8\}\) & \(\{8,4\}\!\ast 64b\) & \(\{4,8\}\!\ast 128b\) & 1024 & 1 \\
86 & \(\{8,4,16\}\) & \(\{8,4\}\!\ast 64a\) & \(\{4,16\}\!\ast 128a\) & 1024 & 1 \\
87 & \(\{8,4,16\}\) & \(\{8,4\}\!\ast 64a\) & \(\{4,16\}\!\ast 128b\) & 1024 & 1 \\
88 & \(\{8,4,16\}\) & \(\{8,4\}\!\ast 64b\) & \(\{4,16\}\!\ast 128a\) & 1024 & 1 \\
89 & \(\{8,4,16\}\) & \(\{8,4\}\!\ast 64b\) & \(\{4,16\}\!\ast 128b\) & 1024 & 1 \\
90 & \(\{8,8,8\}\) & \(\{8,8\}\!\ast 128a\) & \(\{8,8\}\!\ast 128a\) & 1024 & 1 \\
91 & \(\{8,8,8\}\) & \(\{8,8\}\!\ast 128a\) & \(\{8,8\}\!\ast 128b\) & 1024 & 1 \\
92 & \(\{8,8,8\}\) & \(\{8,8\}\!\ast 128a\) & \(\{8,8\}\!\ast 128c\) & 1024 & 1 \\
93 & \(\{8,8,8\}\) & \(\{8,8\}\!\ast 128a\) & \(\{8,8\}\!\ast 128d\) & 1024 & 1 \\
94 & \(\{8,8,8\}\) & \(\{8,8\}\!\ast 128b\) & \(\{8,8\}\!\ast 128a\) & 1024 & 1 \\
95 & \(\{8,8,8\}\) & \(\{8,8\}\!\ast 128b\) & \(\{8,8\}\!\ast 128b\) & 1024 & 1 \\
96 & \(\{8,8,8\}\) & \(\{8,8\}\!\ast 128b\) & \(\{8,8\}\!\ast 128d\) & 1024 & 1 \\
97 & \(\{8,8,8\}\) & \(\{8,8\}\!\ast 128c\) & \(\{8,8\}\!\ast 128a\) & 1024 & 1 \\
98 & \(\{8,8,8\}\) & \(\{8,8\}\!\ast 128c\) & \(\{8,8\}\!\ast 128d\) & 1024 & 1 \\
99 & \(\{8,8,8\}\) & \(\{8,8\}\!\ast 128d\) & \(\{8,8\}\!\ast 128d\) & 1024 & 1 \\
100 & \(\{4,4,4,4\}\) & \(\{4,4,4\}\!\ast 128\) & \(\{4,4,4\}\!\ast 256a\) & 1024 & 1 \\
101 & \(\{4,4,4,4\}\) & \(\{4,4,4\}\!\ast 256a\) & \(\{4,4,4\}\!\ast 256b\) & 1024 & 1 \\
102 & \(\{4,4,4,8\}\) & \(\{4,4,4\}\!\ast 128\) & \(\{4,4,8\}\!\ast 256a\) & 1024 & 1 \\
103 & \(\{4,4,4,8\}\) & \(\{4,4,4\}\!\ast 128\) & \(\{4,4,8\}\!\ast 256b\) & 1024 & 1 \\
104 & \(\{4,4,8,4\}\) & \(\{4,4,8\}\!\ast 256a\) & \(\{4,8,4\}\!\ast 256b\) & 1024 & 1 \\
105 & \(\{4,4,8,4\}\) & \(\{4,4,8\}\!\ast 256a\) & \(\{4,8,4\}\!\ast 256d\) & 1024 & 1 \\
106 & \(\{4,4,8,4\}\) & \(\{4,4,8\}\!\ast 256b\) & \(\{4,8,4\}\!\ast 256a\) & 1024 & 1 \\
107 & \(\{4,4,8,4\}\) & \(\{4,4,8\}\!\ast 256b\) & \(\{4,8,4\}\!\ast 256c\) & 1024 & 1 \\
\end{longtable}
\normalsize

\section{Proof obligations and attribution}
The complete classification and named source results remain attributed
to Gomi, Loyola and De Las Pe\~nas. The parametrization by universal kernels gives a finite exhaustive
search with the marking retained, without an abstract-group catalogue. Central descent and the complete binary lift
argument explain how the source presentations fit that search. The
literal source corrections are part of the declared contract, not
claims of correction priority. Atlas records define named marked
presentations only; their completeness is never assumed.

The exact computational assumptions are GAP's finite permutation and
polycyclic group algorithms, subgroup intersection and completed
Todd--Coxeter enumeration, together with NQ's universal nilpotent
quotient algorithm. A finite image alone would not prove the universal
orders. The one infinite group has the explicitly retained marked
decomposition; it is not certified by a nonterminating enumeration.

The common theory's Lean files verify specified abstract kernel and
intersection steps. They do not prove the full forest quotient theorem,
the class-five census, the 214 raw presentation calculations or the
whole source classification. All those written and computational
obligations remain present. There is no new source credit, question
answer, priority assertion or measured compression factor.

\begin{thebibliography}{9}
\bibitem{UA} Human Compression Project, \emph{Supported kernels and
finite two-group quotients}, companion common paper, this volume.
\bibitem{GLP} Y. Gomi, M. L. Loyola and M. L. A. N. De Las Pe\~nas,
\emph{String C-groups of order 1024}, Contributions to Discrete
Mathematics 13(1) (2018), 1--22, DOI:10.55016/ojs/cdm.v13i1.62475.
\bibitem{MS} P. McMullen and E. Schulte, \emph{Abstract Regular
Polytopes}, Cambridge University Press, 2002.
\end{thebibliography}

\end{document}
